\documentclass{article}

% Encoding and font handling
\usepackage{iftex}
\ifPDFTeX
  \usepackage[T1]{fontenc}
  \usepackage[utf8]{inputenc}
  \usepackage{textcomp}
  \usepackage{lmodern}
\else
  \usepackage{unicode-math}
  \defaultfontfeatures{Scale=MatchLowercase}
  \defaultfontfeatures[\rmfamily]{Ligatures=TeX,Scale=1}
\fi

% Math and color
\usepackage{amsmath, amssymb}
\usepackage{xcolor}

% Bibliography
\usepackage[
  backend=biber,
  style=authoryear,
  natbib=true
]{biblatex}
\addbibresource{literature.bib}

% Paragraph spacing and indentation
\usepackage{parskip} % if available, otherwise fall back to manual
\setlength{\parindent}{0pt}
\setlength{\parskip}{6pt plus 2pt minus 1pt}

% Tables and layout
\usepackage{longtable, booktabs, array, calc}
\usepackage{etoolbox}
\makeatletter
\patchcmd\longtable{\par}{\if@noskipsec\mbox{}\fi\par}{}{}
\makeatother

% Footnotes in tables
\IfFileExists{footnotehyper.sty}{\usepackage{footnotehyper}}{\usepackage{footnote}}
\makesavenoteenv{longtable}

% Links and URLs
\PassOptionsToPackage{unicode}{hyperref}
\PassOptionsToPackage{hyphens}{url}
\usepackage{bookmark}
\IfFileExists{xurl.sty}{\usepackage{xurl}}{}
\urlstyle{same}
\usepackage{hyperref}
\hypersetup{
  hidelinks,
  pdfcreator={LaTeX via custom preamble}
}

% Microtype if available
\IfFileExists{microtype.sty}{
  \usepackage[]{microtype}
  \UseMicrotypeSet[protrusion]{basicmath}
}{}

% Misc
\setcounter{secnumdepth}{-\maxdimen} % remove section numbering
\IfFileExists{upquote.sty}{\usepackage{upquote}}{}
\setlength{\emergencystretch}{3em}
\providecommand{\tightlist}{%
  \setlength{\itemsep}{0pt}\setlength{\parskip}{0pt}}

\author{}
\date{}

\begin{document}



\section{Simulation}\label{simulation}

We start with generating the residuals after random effects/fixed
effects (RE/FE)



  

\begin{equation}\mathbf{\varepsilon}_{RE}\sim N(1,0)
\label{eq:auto_1}
\end{equation}


From this, given GARCH parameters, we can generate volatilities

\[\mathbf{\sigma}^{2} = \exp\left( G\mathbf{+}\mathbf{GA}_{\sigma}^{- 1}\mathbf{AR}_{\sigma}h\left( \mathbf{\varepsilon}_{RE},c \right) \right)\]

And we generate residuals before (RE/FE)



  

\begin{equation}\mathbf{\varepsilon} = \mathbf{\varepsilon}_{RE} + \mathbf{v}_{i} + \mathbf{w}_{t}
\label{eq:auto_2}
\end{equation}


where



  

\begin{equation}\mathbf{v}_{i}\sim N(1,0), \mathbf{w}_{t}\sim N(1,0)
\label{eq:auto_3}
\end{equation}


From this, the observed residuals $\mathbf{u}$ are constructed, given
an ARIMA structure:



  

\begin{equation}
\mathbf{u = A}\mathbf{R}^{- 1}\mathbf{MA \cdot}\mathbf{\varepsilon}
\label{eq:68}
\end{equation}


Exogenous variables $\mathbf{X}$ are then generated, where the first
column $\mathbf{X}_{0}\mathbf{= 1}$ and the second column

\[\mathbf{X}_{1}\mathbf{\sim}N(1,0)\mathbf{\ }\]

Assume a normal distribution:



  

\begin{equation}
f\left( \epsilon_{i,t},\sigma \right) = \frac{1}{\sqrt{2\pi\sigma_{i,t}^{2}}}\exp\left( - \frac{1}{2}\frac{\varepsilon_{i,t}}{\sigma_{i,t}^{2}} \right)
\label{eq:69}
\end{equation}


Where $\varepsilon_{i,t}\ $are residuals after ARIMA and Random
Effects adjustments and $\exp( \cdot )$ is the exponential function.
$\sigma_{i,t}^{2}$ is the EGARCH variance process. The raw
non-stationary residuals are



  

\begin{equation}
\mathbf{u} = \mathbf{y} - \mathbf{X\beta}
\label{eq:70}
\end{equation}


Where $\mathbf{u}$ is a matrix of the ARIMA process $u_{i,t}$,
with dimensions $T \times N$ for a maximum of $T$ dates and $N$
groups. $\mathbf{y}$ is the dependent variable, $\mathbf{X}$ the
matrix of independents and $\mathbf{X\beta}$ the predicted vector.
$\mathbf{X}$ is a three dimensional matrix with dimension
$T \times N \times k$ for $k$ independent variables.

The process is ARMA stationary by the operation:



  

\begin{equation}
\mathbf{\varepsilon} = \mathbf{MA}^{- 1} \cdot \mathbf{AR \cdot u}
\label{eq:71}
\end{equation}


where $\mathbf{\varepsilon}$ is also a $T \times N$ matrix. It is
also allowed that the fixed and random effects occur on the residuals
\emph{before} ARMA transformation, but the consequences of this will be
discussed later.

$\mathbf{\varepsilon}$ may have have fixed/random effects across time
and group:



  

\begin{equation}\mathbf{\varepsilon} = \mathbf{\epsilon} + \mathbf{v}_{i} + \mathbf{w}_{t}
\label{eq:auto_4}
\end{equation}


Where $\mathbf{v}_{i}$ and $\mathbf{w}_{t}$ are fixed and random
effects vectors across group and time respectively, with dimensions
$T \times N$. The residuals before fixed/random effects are then

Residuals after fixed/random effects $\epsilon_{RE,i,t}$ have
time-varying variance that follows an EGARCH process. The EGARCH process
is a function of residuals before random effects:



  

\begin{equation}\mathbf{\sigma}^{2} = \exp\left( G\mathbf{+}\mathbf{GA}_{\sigma}^{- 1}\mathbf{AR}_{\sigma}h\left( \mathbf{\varepsilon}_{RE},c \right) \right)
\label{eq:auto_5}
\end{equation}


\section{Old Simulation}\label{old-simulation}

The random state variables in this system are

\begin{itemize}
\item
  the independent variables
\item
  the residuals
\end{itemize}

Hence a simulation will start by drawing these variables from a normal
distribution.

\subsection{Single equation
simulation:}\label{single-equation-simulation}

It is well known that residual in time series data are often not normal.
In particular this is prevalent in financial and price data, where the
actual distribution often has ``fat tails'', or kurtosis as it is
called. One possible explanation that can generate excess kurtosis is
that the variance is not constant. GARCH is an attempt to model exactly
this. In the simulation function, it will be assumed that it is in deed
the GARCH process that generates kurtosis.

Therefore the standardized residuals
$RE\left( \mathbf{\varepsilon}_{i,t} \right)/\mathbf{\sigma}_{\mathbf{i,t}}$
for firm \emph{i} in the log likelihood function \textbf{Error!
Reference source not found.} will be assumed to be normal, and we define
these as the residual random state variables. Hence, the residuals for
all groups will be generated by drawing each element of a \emph{M} x
\emph{1} vector $\check{\mathbf{\varepsilon}}$ from the standard
normal distribution. We split this vector into \emph{N} groups, where
${\check{\mathbf{\varepsilon}}}_{i}$ represent the standard normal
draws for group \emph{i}.

These standard normal residuals then generates \emph{N} sets of normally
distributed residuals with \emph{random effects}. For each group
\emph{i}, the non standard random effect residual vector is then



  

$\mathbf{\varepsilon}_{\mathbf{i,t}}\mathbf{=}\left( {\check{\mathbf{\varepsilon}}}_{i,t}\mathbf{\sigma}_{i,t}\mathbf{+}\mu_{i}\mathbf{+}\gamma_{t} \right)$ ((74) \\


where $\mu_{i}$ and $\gamma_{t}$ are random effects drawn from a
standard normal distribution for each group and $\mathbf{\sigma}_{i}$
is generated by a E-GARCH process. Hence, the generated standard normal
residuals ${\check{\mathbf{\varepsilon}}}_{i}$ for group \emph{i} are



  

${\check{\mathbf{\varepsilon}}}_{i} = \frac{RE\left( \mathbf{\varepsilon}_{\mathbf{i}} \right)}{\mathbf{\sigma}_{i}} = \frac{\left( \mathbf{\varepsilon}_{\mathbf{i}}\mathbf{-}u_{i} \right)}{\mathbf{\sigma}_{i}}$ ((75) \\


consistent with the model \textbf{Error! Reference source not found.}.
The problem is that the variance process $\mathbf{\sigma}_{i}$ in (74)
depends on $\mathbf{\varepsilon}_{\mathbf{i}}$, which again depends on
the variance process $\mathbf{\sigma}_{i}$.

Let us start with the variance process. It follows from \textbf{Error!
Reference source not found.}, (20) and (24) in the E-GARCH model
(section 4) that



  

\begin{equation}
\ln\left( \mathbf{\sigma}_{i}^{2} \right) = \mathbf{W}_{i}\mathbf{\omega}^{\mathbf{'}}\mathbf{+}\left( \mu\frac{1}{T_{i}}\mathbf{1}\mathbf{1}^{\mathbf{'}}\mathbf{+}{AR}_{\sigma}^{- 1}MA_{\sigma} \right)\ln\left( \mathbf{\varepsilon}_{i}^{2} + c \right)
\label{eq:76}
\end{equation}


Inserting (74) yields



  

\begin{equation}
\ln\left( \mathbf{\sigma}_{i}^{2} \right) = \mathbf{W}_{i}\mathbf{\omega}^{\mathbf{'}}\mathbf{+}\left( \mu\frac{1}{T_{i}}\mathbf{1}\mathbf{1}^{\mathbf{'}}\mathbf{+}{AR}_{\sigma}^{- 1}MA_{\sigma} \right)\ln\left( \left( {\check{\mathbf{\varepsilon}}}_{i}\mathbf{\sigma}_{i}\mathbf{+}u_{i} \right)^{2} + c \right)
\label{eq:77}
\end{equation}


It is not possible to solve this with respect to
$\mathbf{\sigma}_{i}$, unless we do some simplifying assumtions.
$\ln\left( \left( {\check{\mathbf{\varepsilon}}}_{i}\mathbf{\sigma}_{i}\mathbf{+}u_{i} \right)^{2} + c \right)$

Since $\mathbf{\varepsilon}_{i}^{2}$ are the residuals resulting from
applying the ARIMA filter (section 3) on observed residuals

Define the standardized residuals

\[{\overline{\mathbf{\varepsilon}}}_{j,i}^{2} = \frac{RE\left( \mathbf{\varepsilon}_{i}^{2} \right)}{\mathbf{\sigma}_{i}^{2}}\]

I will assume that these residuals are random noise with standard normal
distribution. As before $RE\left( \mathbf{\cdot} \right)$ represents
the random effects transformation of the residuals. Hence, we can

\[RE\left( \mathbf{\varepsilon}_{i}^{2} \right) = {\overline{\mathbf{\varepsilon}}}_{j,i}^{2}\mathbf{\sigma}_{i}^{2}\]

We now need a simplifying assumption in order to write
$\mathbf{\sigma}_{i}^{2}$ as a function of the state variables
${\overline{\mathbf{\varepsilon}}}_{j,i}^{2}$. The mathematically
correct way to invert the random effects function
$RE\left( \mathbf{\varepsilon}_{i}^{2} \right)$, would be $
$\[RE\left( \mathbf{\varepsilon}_{i}^{2} \right) = {\overline{\mathbf{\varepsilon}}}_{j,i}^{2}\mathbf{\sigma}_{i}^{2}\]

\[RE\left( \mathbf{\varepsilon}_{i}^{2} \right) = {\overline{\mathbf{\varepsilon}}}_{j,i}^{2}\mathbf{\sigma}_{i}^{2}\]

\[\ln\left( \mathbf{\sigma}_{i}^{2} \right) = \mathbf{W}_{i}\mathbf{\omega}^{\mathbf{'}}\mathbf{+}\left( \mu\frac{1}{T_{i}}\mathbf{1}\mathbf{1}^{\mathbf{'}}\mathbf{+}{AR}_{\sigma}^{- 1}MA_{\sigma} \right)\ln\left( {\overline{\mathbf{\varepsilon}}}_{j,i}^{2}\mathbf{\sigma}_{i}^{2} + c \right)\]

where ${\overline{\mathbf{\varepsilon}}}_{j,i}^{2}$ are the
standardized residuals. In order to solve this equation, we need to
assume \emph{c}=0:

\[\ln\left( \mathbf{\sigma}_{i}^{2} \right) = \mathbf{W}_{i}\mathbf{\omega}^{\mathbf{'}}\mathbf{+}\left( \mu\frac{1}{T_{i}}\mathbf{1}\mathbf{1}^{\mathbf{'}}\mathbf{+}{AR}_{\sigma}^{- 1}MA_{\sigma} \right)\left( \ln\left( {RE}^{- 1}\left( {\overline{\mathbf{\varepsilon}}}_{j,i}^{2} \right) \right) + \ln\left( \mathbf{\sigma}_{i}^{2} \right) \right)\]

Solving for $\ln\left( \mathbf{\sigma}_{i}^{2} \right)$:

\[\ln\left( \mathbf{\sigma}_{i}^{2} \right) = \left( I - \left( \mu\frac{1}{T_{i}}\mathbf{1}\mathbf{1}^{\mathbf{'}}\mathbf{+}{AR}_{\sigma}^{- 1}MA_{\sigma} \right) \right)^{- 1}\left( \mathbf{W}_{i}\mathbf{\omega}^{\mathbf{'}}\mathbf{+}\left( \mu\frac{1}{T_{i}}\mathbf{1}\mathbf{1}^{\mathbf{'}}\mathbf{+}{AR}_{\sigma}^{- 1}MA_{\sigma} \right) \right)\ln\left( {RE}^{- 1}\left( {\overline{\mathbf{\varepsilon}}}_{j,i}^{2} \right) \right)\]

\subsection{The independent variables}\label{the-independent-variables}

The independent variables will drawn given a custom covariance matrix
$\mathbf{\Sigma}$\textbf{,} by decomposing it as
$\mathbf{\Sigma} = \mathbf{P'P}$. Say $\mathbf{\Sigma}$ is a
\emph{k} x \emph{k} matrix and the total size of the sample is \emph{M}.
If we makes \emph{k}⋅\emph{n} draws from the standard normal
distribution and place them in the \emph{M} x \emph{k} matrix
$\mathbf{Z}$, then we can generate a set of independent variables with
the desired covariance structure and unity mean as
$\mathbf{X} = \mathbf{P'}\mathbf{Z}\mathbf{+ 1}$, where \textbf{1} is
a \emph{k} x \emph{1} vector of ones. Of course, one could easily allow
means to be customizable, but there is difficult to see a compelling
reason for that. The same holds for the coefficients of the independent
variables. These are also assumed to be one in the simulation, that is
$\mathbf{\beta = 1}$.
\printbibliography
\end{document}
